% Part: many-valued-logic
% Chapter: sequent-calculus
% Section: rules-and-proofs

\documentclass[../../../include/open-logic-section]{subfiles}

\begin{document}

\olfileid{mvl}{seq}{rul}

\olsection{বিধি ও \usetoken{P}{derivation}}

পরের আলোচনায় $\Gamma, \Delta, \Pi, \Lambda$ দিয়ে
!!{sentence}s-এর সসীম অনুক্রম বোঝানো হবে।

\begin{defn}[সিকোয়েন্ট]
কোনো \emph{$n$-পাশবিশিষ্ট সিকোয়েন্ট} হলো নিচের রূপের অভিব্যক্তি:
\[
\Gamma_1 \nSequent \dots \nSequent \Gamma_n
\]
এখানে প্রতিটি $\Gamma_i$ ভাষা~$\Lang L$-এর !!{sentence}s-এর
একটি সসীম (সম্ভবত ফাঁকা) অনুক্রম।
% BN-SRC-331 TARGET-CORRECTION-BEGIN
\par\smallskip\noindent\textnormal{\small\textbf{উৎস-সংশোধন (BN-SRC-331):} হিমায়িত উৎসে ``প্রতিটি $\Gamma_1$'' বলা হয়েছে। উপরের সাধারণ $n$-পাশবিশিষ্ট সিকোয়েন্টে প্রতিটি অবস্থানের অনুক্রম বোঝাতে এখানে $\Gamma_i$ ব্যবহার করা হয়েছে।} % BN-SRC-331 TARGET-CORRECTION-END
\end{defn}

\begin{defn}[প্রারম্ভিক সিকোয়েন্ট]
কোনো \emph{$n$-পাশবিশিষ্ট প্রারম্ভিক সিকোয়েন্ট} হলো
$!A \nSequent \dots \nSequent !A$ রূপের একটি $n$-পাশবিশিষ্ট
সিকোয়েন্ট, যেখানে $!A$ ভাষাটির যেকোনো !!{sentence}।

ভাষাটিতে যদি কোনো $0$-স্থানবিশিষ্ট সংযোজক~$\star$ থাকে,
অর্থাৎ কোনো বচনধ্রুবক থাকে, তবে $\dots \nSequent
\star \nSequent \dots$ সিকোয়েন্টটিকেও প্রারম্ভিক ধরি।
এখানে $\star$ থাকে $\tf{\star} \in V$ সত্যমানের সঙ্গে
সংশ্লিষ্ট অবস্থানে; অন্য অবস্থানগুলি ফাঁকা থাকে।
\end{defn}

$n$-মানী যুক্তিবিদ্যা~$\Log{L}$-এর প্রতিটি সংযোজকের জন্য,
$\Log{L}$-এ সংযোজকটি যতগুলি সত্যমান নিতে পারে, প্রতিটির
জন্য একটি করে যৌক্তিক বিধি আছে। $\Log{L}$-এর $n$-পাশবিশিষ্ট
সিকোয়েন্ট কলনে !!^{derivation}s হলো সিকোয়েন্টের বৃক্ষ।
সর্বোচ্চ সিকোয়েন্টগুলি প্রারম্ভিক সিকোয়েন্ট; আর কোনো সিকোয়েন্ট
এক বা একাধিক অন্য সিকোয়েন্টের নিচে থাকলে সেটি অবশ্যই
$\Log{L}$-এর সংযোজকগুলির কোনো অনুমান-বিধি সঠিকভাবে প্রয়োগ করে
পাওয়া যেতে হবে।

\begin{defn}[উপপাদ্য]
কোনো !!^a{sentence}~$!A$ $n$-মানী যুক্তিবিদ্যা~$\Log{L}$-এর
\emph{উপপাদ্য} তখনই, যখন $\Log{L}$-এর প্রতিটি মনোনীত সত্যমানের
সংশ্লিষ্ট অবস্থানে $!A$-যুক্ত $n$-সিকোয়েন্টটির একটি
!!a{derivation} থাকে। $!A$ উপপাদ্য হলে লিখি
$\Proves[\Log{L}] !A$, আর উপপাদ্য না হলে
$\Proves/[\Log{L}] !A$।
\end{defn}

\begin{defn}[!!^{derivability}]
কোনো !!^a{sentence}~$!A$ $n$-মানী যুক্তিবিদ্যা~$\Log{L}$-এ
!!{sentence}s-এর সেট~$\Gamma$ থেকে \emph{!!{derivable}}, অর্থাৎ
$\Gamma \Proves[\Log{L}] !A$, তখনই, যখন একটি সসীম
উপসেট~$\Gamma_0 \subseteq \Gamma$ এবং $\Gamma_0$-র
!!{sentence}s-এর একটি অনুক্রম $\Gamma_0'$ আছে, যাতে নিচের
সিকোয়েন্টটির একটি !!a{derivation} থাকে:
\[ \Lambda_1 \nSequent \dots \nSequent \Lambda_n \]
এখানে $i$-তম অবস্থান মনোনীত সত্যমানের হলে $\Lambda_i$ হলো
$!A$, অন্যথায় $\Gamma_0'$। $!A$ যদি $\Gamma$ থেকে
!!{derivable} না হয়, তবে লিখি $\Gamma \Proves/ !A$।
\end{defn}

উদাহরণ হিসেবে, $3$-মানী \L ukasiewicz যুক্তিবিদ্যার একটি
$3$-পাশবিশিষ্ট সিকোয়েন্ট কলন আছে। $3$-পাশবিশিষ্ট
$\Gamma \nSequent \Pi \nSequent \Delta$ সিকোয়েন্টে
$\Gamma$-র অবস্থান~$\False$, $\Delta$-র অবস্থান~$\True$,
আর $\Pi$-র অবস্থান~$\Undef$-এর সঙ্গে সংশ্লিষ্ট।
প্রারম্ভিক সিকোয়েন্টগুলি $!A \nSequent !A \nSequent !A$।
শুধু $\True$ মনোনীত বলে $\Gamma \Proves[\LogLuk[3]] !A$
তখনই, যখন $\Gamma \nSequent \Gamma \nSequent !A$
সিকোয়েন্টটির একটি !!a{derivation} থাকে। ($\Undef$-ও
মনোনীত হলে $\Gamma \nSequent !A \nSequent !A$-এর একটি
!!a{derivation} দরকার হতো।)

\end{document}
